\section{MCP: Model Context Protocol}

\eng{Model Context Protocol} (\eng{MCP}) הוא תקן שמגדיר כיצד סוכני \eng{AI}
מתחברים לשירותים חיצוניים --- מסדי נתונים, מאגרי קבצים, כלי פיתוח ועוד
\cite{anthropic2024mcp}. ההרצאה מציגה את \eng{MCP} כאלטרנטיבה מודרנית
לעבודה ישירה עם \eng{API} ו-\eng{Skills} סטטיים.

% MCP vs Skill vs API
\begin{figure}[H]
\centering
\begin{tikzpicture}[
  box/.style={draw, rounded corners=3pt, minimum width=2.5cm, minimum height=1.6cm,
    align=center, font=\footnotesize},
  arr/.style={-{Stealth[length=2mm]}, thick}
]
  \node[box, fill=accentblue!12, draw=accentblue] (agent) at (0,0) {\eng{AI Agent}};
  \node[box, fill=lightgray, draw=secondary] (skill) at (-4,-2.2) {\eng{Skill}\\סטטי מקומי};
  \node[box, fill=accentamber!12, draw=accentamber] (api) at (0,-2.2) {\eng{API}\\נקודת קצה};
  \node[box, fill=accentgreen!12, draw=accentgreen] (mcp) at (4,-2.2) {\eng{MCP Server}\\מחובר מרוחק};
  \node[box, fill=primary!8, draw=primary] (sdk) at (0,-4.2) {\eng{SDK / Software}};
  \draw[arr, accentblue] (agent) -- (skill);
  \draw[arr, accentamber] (agent) -- (api);
  \draw[arr, accentgreen] (agent) -- (mcp);
  \draw[arr, secondary] (skill) -- (sdk);
  \draw[arr, secondary] (api) -- (sdk);
  \draw[arr, secondary] (mcp) -- (sdk);
\end{tikzpicture}
\caption{השוואה מושגית: סוכן יכול לגשת ליכולות תוכנה דרך \eng{Skill}, \eng{API} או \eng{MCP}}
\label{fig:mcp-comparison}
\end{figure}


\subsection{מהו MCP?}

\eng{MCP} הוא ``שרת'' שרץ בנקודת קצה --- במחשב המקומי או אצל ספק התוכנה ---
ומציע לסוכן רשימת כלים (\eng{tools}), משאבים (\eng{resources}) ותבניות
פרומפט. הסוכן ``מכיר'' את ה-\eng{MCP} הרשום אצלו (בדומה לרשימת \eng{Skills})
ומפעיל אותו כשהמשימה מתאימה.

\subsection{MCP מול API ישיר}

בעבודה מול \eng{API} ישיר, המפתח חייב לדעת כתובת, אימות, סכמת בקשות
ועדכוני גרסה. ב-\eng{MCP}, ספק התוכנה מטפל בשכבה זו; הסוכן מגלה יכולות
דינמית מהשרת.

\begin{insightbox}
ההבדל המרכזי מההרצאה: \eng{Skill} סטטי דורש התקנה מחדש כשה-\eng{API}
משתנה; \eng{MCP} מחובר ``אונליין'' ומתעדכן אוטומטית מול היצרן.
\end{insightbox}

\subsection{מודל שרת-לקוח}

\eng{MCP} מבוסס על מודל שרת ולקוח, בדומה לגלישה באינטרנט: הסוכן ``גולש''
לשרת ומבצע פעולות מורשות. אין צורך לדעת כל פרט טכני על מיקום התוכנה ---
רק את נקודת הקצה של ה-\eng{MCP}.

\subsection{דוגמת תצורה (מושגית)}

\begin{lstlisting}[caption={דוגמת הגדרת MCP בקובץ JSON}]
{
  "mcpServers": {
    "filesystem": {
      "command": "npx",
      "args": ["-y", "@modelcontextprotocol/server-filesystem", "./docs"]
    },
    "browser": {
      "command": "cursor-ide-browser",
      "args": []
    }
  }
}
\end{lstlisting}

\subsection{MCP בתוך ארכיטקטורת הסוכן}

כפי שהוסבר: הסוכן מכיל פנייה ל-\eng{MCP}. כמו חיפוש \eng{Skills}, הוא
בוחר את השרת המתאים. בתוך ה-\eng{MCP} עדיין קיים \eng{API} שמדבר עם
ה-\eng{SDK} --- אך הסוכן לא חשוף לפרטים אלה.

\agenttable{
\caption{השוואה: \eng{Skill}+\eng{API} מול \eng{MCP}}
\label{tab:mcp-vs-skill}
\begin{tabular}{@{}p{3cm}p{4.5cm}p{4.5cm}@{}}
\toprule
\textbf{מאפיין} & \textbf{Skill + API} & \textbf{MCP} \\
\midrule
מיקום & מקומי, סטטי & שרת מרוחק/מקומי \\
עדכון & ידני & אוטומטי \\
גילוי יכולות & מתוך קובץ Skill & מתוך שרת \\
תקשורת בין סוכנים & שפה חופשית & שפה חופשית \\
\bottomrule
\end{tabular}
}

\subsection{מגבלות ואבטחה}

\begin{itemize}[rightmargin=2em]
  \item הרשאות --- אילו כלי \eng{MCP} מורשים לסוכן?
  \item אימות --- מפתחות \eng{API} וסודות לא בפרומפט.
  \item זמינות --- תלות ברשת ובשרת חיצוני.
  \item תאימות --- לא כל מוצר תומך ב-\eng{MCP} עדיין.
\end{itemize}

\begin{warningbox}
חשיפת \eng{MCP} עם גישה לקבצים מערכת או דפדפן דורשת הגדרת הרשאות
זהירה --- סוכן שגוי עלול למחוק קבצים או לשלוח מידע רגיש.
\end{warningbox}

\subsection{קשר לפרויקט הנוכחי}

ב-\eng{Cursor}, שרתי \eng{MCP} כמו \eng{cursor-ide-browser} ו-\eng{playwright}
מאפשרים לסוכן לפעול בסביבת הפיתוח. מטלה 04 עצמה מדגימה עבודה טקסטואלית
(\LaTeX) --- גישה עדיפה על סימולציית \eng{GUI} לעריכת מסמכים.
